\documentclass[12pt,a4paper]{ctexart}

\usepackage{amsmath,amssymb}
\usepackage{geometry}
\usepackage{enumitem}
\usepackage{hyperref}
\usepackage{booktabs}
\usepackage{listings}

\geometry{left=2.5cm,right=2.5cm,top=2.5cm,bottom=2.5cm}
\setlist{nosep,leftmargin=2em}

\title{第1章\quad 质点运动学}
\author{}
\date{}

\begin{document}

\maketitle

\section{知识点总结}

\subsection{基本概念}

\begin{itemize}
    \item \textbf{质点}：有质量而无形状和大小的几何点，是理想模型。可视为质点的两种情况：\textcircled{1} 两相互作用物体间距离远大于本身线度；\textcircled{2} 物体作平动。
    \item \textbf{质点系}：若干质点的集合。
    \item \textbf{参照物 / 参考系}：参照物 + 坐标系 + 时钟。运动学中参考系可任选；不同参考系下物体的运动状态和轨迹可能不同；同一参考系下物体运动状态与坐标系选择无关，只是数学描述形式不同。
    \item \textbf{常见坐标系}：直角坐标系、球坐标系、极坐标系、自然坐标系。
\end{itemize}

\subsection{质点位置的确定}

\begin{itemize}
    \item \textbf{直角坐标法}：用 $P(x, y, z)$ 三个坐标确定位置。
    \item \textbf{位矢法}：$\vec{r} = x\vec{i} + y\vec{j} + z\vec{k}$，大小 $|\vec{r}| = \sqrt{x^2 + y^2 + z^2}$，方向用方向余弦 $\cos\alpha, \cos\beta, \cos\gamma$ 表示。
    \item \textbf{自然坐标法}：$s = f(t)$，用于已知运动轨迹的情况。
\end{itemize}

\subsection{运动学方程}

\begin{itemize}
    \item 直角坐标下：$\vec{r} = \vec{r}(t) = x(t)\vec{i} + y(t)\vec{j} + z(t)\vec{k}$，各基矢量方向不变。
    \item 自然坐标下：$s = f(t)$。
    \item 已知运动学方程可求质点的运动轨迹、速度和加速度。
\end{itemize}

\subsection{位移、速度、加速度}

\begin{itemize}
    \item \textbf{位移}：$\Delta\vec{r} = \vec{r}(t+\Delta t) - \vec{r}(t)$，是矢量，$|\Delta\vec{r}| \le \Delta s$（路程）。
    \item \textbf{平均速度}：$\langle\vec{v}\rangle = \dfrac{\Delta\vec{r}}{\Delta t}$。
    \item \textbf{瞬时速度}：$\vec{v} = \dfrac{\mathrm{d}\vec{r}}{\mathrm{d}t}$，方向沿轨道切向指向质点前进方向。
    \item \textbf{瞬时速率}：$v = \dfrac{\mathrm{d}s}{\mathrm{d}t} = |\vec{v}|$，注意 $v \ne \dfrac{\mathrm{d}r}{\mathrm{d}t}$。
    \item \textbf{平均加速度}：$\langle\vec{a}\rangle = \dfrac{\Delta\vec{v}}{\Delta t}$。
    \item \textbf{瞬时加速度}：$\vec{a} = \dfrac{\mathrm{d}\vec{v}}{\mathrm{d}t} = \dfrac{\mathrm{d}^2\vec{r}}{\mathrm{d}t^2}$。
    \item 加速度反映速度方向和量值的变化；其法向分量总是指向轨迹曲线凹的一面。
    \item 加速度与速度夹角判断运动性质：$0^\circ/180^\circ$ 直线运动，$90^\circ$（匀速率）圆周运动，$<90^\circ$ 加速，$>90^\circ$ 减速。
\end{itemize}

\subsection{直角坐标表示}

\begin{itemize}
    \item 速度分量：$v_x = \dfrac{\mathrm{d}x}{\mathrm{d}t},\; v_y = \dfrac{\mathrm{d}y}{\mathrm{d}t},\; v_z = \dfrac{\mathrm{d}z}{\mathrm{d}t}$。
    \item 速度大小：$v = \sqrt{v_x^2 + v_y^2 + v_z^2}$。
    \item 加速度分量：$a_x = \dfrac{\mathrm{d}^2 x}{\mathrm{d}t^2},\; a_y = \dfrac{\mathrm{d}^2 y}{\mathrm{d}t^2},\; a_z = \dfrac{\mathrm{d}^2 z}{\mathrm{d}t^2}$。
    \item 加速度大小：$|\vec{a}| = \sqrt{a_x^2 + a_y^2 + a_z^2}$。
\end{itemize}

\subsection{自然坐标表示}

\begin{itemize}
    \item 速度：$\vec{v} = \dfrac{\mathrm{d}s}{\mathrm{d}t}\vec{\tau}$（沿切向）。
    \item \textbf{切向加速度}：$a_\tau = \dfrac{\mathrm{d}v}{\mathrm{d}t}$，反映速度大小变化的快慢。
    \item \textbf{法向加速度}：$a_n = \dfrac{v^2}{\rho}$，$\rho$ 为曲率半径，反映速度方向变化的快慢。
    \item 总加速度：$\vec{a} = a_\tau\vec{\tau} + a_n\vec{n}$，大小 $a = \sqrt{a_\tau^2 + a_n^2}$，夹角 $\tan\theta = a_n/a_\tau$。
    \item 直线运动和圆周运动是平面曲线运动的特例。
\end{itemize}

\subsection{圆周运动的角量描述}

\begin{itemize}
    \item \textbf{角位置}：$\theta = \theta(t)$；\textbf{角位移}：$\Delta\theta$，按右手法则，逆时针为正。
    \item \textbf{角速度}：$\vec{\omega} = \dfrac{\mathrm{d}\theta}{\mathrm{d}t}\vec{k}$，方向由右手法则确定，单位 rad/s。
    \item \textbf{角加速度}：$\vec{\beta} = \dfrac{\mathrm{d}\vec{\omega}}{\mathrm{d}t} = \dfrac{\mathrm{d}^2\theta}{\mathrm{d}t^2}\vec{k}$，单位 $\mathrm{rad}/\mathrm{s}^2$。
\end{itemize}

\subsection{角量与线量关系}

\begin{itemize}
    \item 位移：$|\mathrm{d}\vec{r}| = r\,\mathrm{d}\theta$，矢量式 $\mathrm{d}\vec{r} = \mathrm{d}\theta\,\vec{k}\times\vec{r}$。
    \item $s = R\theta,\; v = R\omega,\; a_\tau = R\beta,\; a_n = R\omega^2 = \dfrac{v^2}{R}$。
    \item 矢量式：$\vec{v} = \vec{\omega}\times\vec{r}$，$\vec{a} = \vec{\beta}\times\vec{r} + \vec{\omega}\times\vec{v}$。
\end{itemize}

\subsection{运动学二类问题}

\begin{itemize}
    \item \textbf{第一类问题}（求导）：已知运动学方程 $\vec{r}(t)$，求 $\vec{v}, \vec{a}$。
    \item \textbf{第二类问题}（积分）：已知加速度 $\vec{a}$ 和初始条件，求 $\vec{v}(t), \vec{r}(t)$。
\end{itemize}

\subsection{参考系变换（伽利略变换）}

\begin{itemize}
    \item 仅适用于参照系间只发生平动、低速宏观机械运动。
    \item 速度变换：$\vec{v}_{\text{绝对}} = \vec{v}_{\text{相对}} + \vec{v}_{\text{牵连}}$。
    \item 加速度变换：$\vec{a}_{\text{绝对}} = \vec{a}_{\text{相对}} + \vec{a}_{\text{牵连}}$。
\end{itemize}

\section{公式汇总}

\begin{enumerate}[leftmargin=3em,label=\textbf{公式\arabic*}：,ref=公式\arabic*]

    \item $\vec{r} = x\vec{i} + y\vec{j} + z\vec{k}$

    说明：位矢在直角坐标系中的表示，$x,y,z$ 为质点坐标，$\vec{i},\vec{j},\vec{k}$ 为单位矢量。

    \item $\Delta\vec{r} = \vec{r}(t+\Delta t) - \vec{r}(t)$

    说明：位移矢量，$\Delta t$ 时间内位矢的增量，$|\Delta\vec{r}|$ 是位移大小（标量），与路程 $\Delta s$ 不同。

    \item $\vec{v} = \dfrac{\mathrm{d}\vec{r}}{\mathrm{d}t}$

    说明：瞬时速度，位矢对时间的一阶导数，方向沿轨迹切向指向质点前进方向。

    \item $v = \dfrac{\mathrm{d}s}{\mathrm{d}t} = |\vec{v}|$

    说明：瞬时速率（标量），等于速度的模；注意 $v \ne \mathrm{d}r/\mathrm{d}t$。

    \item $\vec{a} = \dfrac{\mathrm{d}\vec{v}}{\mathrm{d}t} = \dfrac{\mathrm{d}^2\vec{r}}{\mathrm{d}t^2}$

    说明：瞬时加速度，速度对时间的一阶导数，反映速度方向与大小的变化。

    \item $v_x = \dfrac{\mathrm{d}x}{\mathrm{d}t},\; v_y = \dfrac{\mathrm{d}y}{\mathrm{d}t},\; v_z = \dfrac{\mathrm{d}z}{\mathrm{d}t}$；$a_x = \dfrac{\mathrm{d}^2x}{\mathrm{d}t^2},\; a_y = \dfrac{\mathrm{d}^2y}{\mathrm{d}t^2},\; a_z = \dfrac{\mathrm{d}^2z}{\mathrm{d}t^2}$

    说明：直角坐标系中速度、加速度的分量形式。

    \item $\vec{v} = \dfrac{\mathrm{d}s}{\mathrm{d}t}\vec{\tau}$

    说明：自然坐标系下的速度表示，$\vec{\tau}$ 为切向单位矢量。

    \item $a_\tau = \dfrac{\mathrm{d}v}{\mathrm{d}t},\quad a_n = \dfrac{v^2}{\rho}$

    说明：自然坐标系下的切向、法向加速度。$\rho$ 为曲率半径。

    \item $\vec{a} = \dfrac{\mathrm{d}v}{\mathrm{d}t}\vec{\tau} + \dfrac{v^2}{\rho}\vec{n}$

    说明：总加速度在自然坐标系中的分解表示，大小 $a=\sqrt{a_\tau^2+a_n^2}$。

    \item $\omega = \dfrac{\mathrm{d}\theta}{\mathrm{d}t},\quad \beta = \dfrac{\mathrm{d}\omega}{\mathrm{d}t} = \dfrac{\mathrm{d}^2\theta}{\mathrm{d}t^2}$

    说明：圆周运动的角速度与角加速度，$\omega$ 单位 rad/s，$\beta$ 单位 $\mathrm{rad}/\mathrm{s}^2$。

    \item $s = R\theta,\quad v = R\omega,\quad a_\tau = R\beta,\quad a_n = R\omega^2 = \dfrac{v^2}{R}$

    说明：圆周运动中线量与角量的关系，$R$ 为圆周半径。

    \item $\vec{v} = \vec{\omega}\times\vec{r}$，$\vec{a} = \vec{\beta}\times\vec{r} + \vec{\omega}\times\vec{v}$

    说明：速度、加速度与角速度、角加速度的矢量叉乘关系，方向由右手法则确定。

    \item $\vec{v}_{\text{绝对}} = \vec{v}_{\text{相对}} + \vec{v}_{\text{牵连}}$

    说明：伽利略速度变换定理，仅适用于参照系间只发生平动的情况。

    \item $\vec{a}_{\text{绝对}} = \vec{a}_{\text{相对}} + \vec{a}_{\text{牵连}}$

    说明：伽利略加速度变换定理。

    \item 匀变速直线运动三公式：$v = v_0 + at,\; x = x_0 + v_0 t + \dfrac{1}{2}at^2,\; v^2 = v_0^2 + 2a(x-x_0)$

    说明：匀变速直线运动中速度、时间和位移的关系，可由 $a=\mathrm{d}v/\mathrm{d}t$ 积分得到。

\end{enumerate}

\end{document}
